\section{数值验证与误差分析}
数值验证围绕三个问题展开：离散格式是否正确，空间与时间分辨率是否足够，以及终点结果是否稳定。表\ref{tab:verification-map}列出各项检验的参照对象与作用。

\begin{table}[htbp]\centering\small
\caption{验证项目、参照对象与可识别的错误}\label{tab:verification-map}
\begin{tabularx}{\linewidth}{lXX}
\toprule
检验 & 参照对象 & 主要识别问题\\
\midrule
固定域解析模态 & Robin边界Bessel解与精确越阈时间 & 柱坐标算子、中心和表面符号、事件定位\\
移动域非均匀模态 & 变半径零通量解析解 & 几何因子是否真正随时间更新\\
几何因子对照算例 & 只把$R(t)^{-2}$冻结为$R_0^{-2}$ & 固定几何因子对求解误差的影响\\
水分收支 & 离散均值与累计边界损失 & 面通量、积分权重和干基归一化\\
雅可比方向差分 & 实际右端的中心差分 & 状态及交叉导数是否一致\\
网格与容差复验 & 同一工况下更细离散结果 & 本题场与终点对数值设置的敏感性\\
\bottomrule
\end{tabularx}
\end{table}

\subsection{解析极限、封闭场与导数检查}
对常扩散系数$d_0=10^{-5}$平方米每秒、半径0.02米、$C_e=0.05$、$h_m=d_0/R$的算例，取式\eqref{eq:bessel}的Bessel模态初值。其中$\beta$为$\beta J_1(\beta)=J_0(\beta)$的第一正根，解析解同时满足中心对称和Robin表面条件。该组参数专用于数值验证。

\begin{equation}\label{eq:bessel}
C(r,t)=0.05+0.20J_0\left(\beta\frac rR\right)
       \exp\left(-\frac{d_0\beta^2}{R^2}t\right),\qquad
t_{\rm exact}=\frac{R^2\log2}{d_0\beta^2}.
\end{equation}

采用160个区间时，精确越阈时间为17.581493358秒，数值事件时间为17.581482896秒，相差约$1.05\times10^{-5}$秒。该算例共同检验柱坐标离散、中心处理、表面通量符号和事件定位。另取零传质系数与均匀初值，离散右端为零，封闭均匀场保持不变。

以方向差分比较$J(y)p$和$[F(y+\varepsilon p)-F(y-\varepsilon p)]/(2\varepsilon)$，检查解析雅可比与方程右端的一致性。检验覆盖四问，并在常湿、低湿、均匀及加密网格下比较16组块导数，最大相对差约为$3.78\times10^{-10}$。

\subsection{非均匀初值的移动域解析基准}
为检验变化半径下的扩散速度，取$R(t)=R_0(1-\gamma t)$，其中$R_0=0.02$米、$\gamma=1/7200$秒$^{-1}$、$0\leq t\leq3600$秒，半径由0.02米缩至0.01米。扩散系数固定为$D_0=10^{-8}$平方米每秒，中心和表面均为零通量。取$J_1(\lambda)=0$的第一正根$\lambda=3.8317059702$，式\eqref{eq:moving-bessel}给出非均匀初值及其解析解。

\begin{equation}\label{eq:moving-bessel}
\begin{aligned}
C(x,t)&=1+0.3J_0(\lambda x)
 \exp\left[-\lambda^2D_0 I(t)\right],\\
I(t)&=\int_0^t\frac{\mathrm ds}{R(s)^2}
 =\frac{t}{R_0^2(1-\gamma t)}.
\end{aligned}
\end{equation}

由$J_0'(z)=-J_1(z)$可知，该解在$x=0,1$处均满足零梯度，Bessel特征方程保证其满足常系数水分方程。空间非均匀模态按$I(t)$衰减，能检验几何因子对扩散速度的影响。算例使用本文采用的水分离散算子，仅替换物性、半径轨迹和边界条件。

空间加密时，相对容差取$10^{-11}$，最大时间步取14.0625秒，在0至3600秒的129个等间隔时刻、各网格全部节点上比较数值解和解析解。表\ref{tab:moving-benchmark}给出最大绝对误差。随着区间数增加，误差逐级下降，640区间时为$8.11\times10^{-7}$，低于$10^{-5}$的误差判据。

\begin{table}[htbp]\centering\small
\caption{移动域解析算例的空间加密与几何因子对照}\label{tab:moving-benchmark}
\begin{tabular}{lrr}
\toprule
算子 & 径向区间数 & 最大解析对比误差/(kg/kg)\\
\midrule
随时间更新$R(t)^{-2}$ & 40 & $2.05124\times10^{-4}$\\
随时间更新$R(t)^{-2}$ & 80 & $5.17059\times10^{-5}$\\
随时间更新$R(t)^{-2}$ & 160 & $1.29592\times10^{-5}$\\
随时间更新$R(t)^{-2}$ & 320 & $3.24227\times10^{-6}$\\
随时间更新$R(t)^{-2}$ & 640 & $8.10753\times10^{-7}$\\
故意冻结为$R_0^{-2}$ & 640 & $5.86815\times10^{-2}$\\
\bottomrule
\end{tabular}
\end{table}

时间加密采用160区间均匀网格，相对容差依次为$10^{-3}$、$10^{-5}$、$10^{-7}$，最大时间步依次为450、225、112.5秒。以同一半离散矩阵$L_N$的指数解$\exp[D_0I(t)L_N]C(0)$为参照，时间误差依次为$1.60038\times10^{-3}$、$2.56224\times10^{-5}$、$2.65945\times10^{-7}$，最细设置满足$5\times10^{-6}$的误差判据。该设置与连续Bessel解的总误差为$6.37797\times10^{-6}$，其中包含固定空间网格的离散误差。

对照算例保持相同的移动半径、非均匀初值和零通量边界，仅将离散方程中的$R(t)^{-2}$固定为$R_0^{-2}$。误差增至约$5.87\times10^{-2}$，明显超过$10^{-5}$的判据。该情形仍保持离散水量，表明几何因子的正确性需结合空间分布比较判断。

正确离散下，平均含水率相对初始离散加权均值的漂移不超过$6.66\times10^{-16}$。此处比较同一离散求积下的前后变化，与初值相对连续均值1的求积误差分开。原有固定域、收缩域水分平衡检验及无通量均匀场检验均保持通过。

\subsection{归一化水分平衡}
在空间均匀干固体密度条件下，径向平均为$\overline C=2\int_0^1xC\,dx$。对水分方程积分并代入表面条件，得到式\eqref{eq:balance}。令累计归一化损失$\Lambda$从零开始与温湿状态同步积分，则$\overline C+\Lambda$应保持初值2.55。

\begin{equation}\label{eq:balance}
\frac{d\overline C}{dt}=-\frac{2h_m}{R(t)}(C_s-C_e),\qquad
\dot\Lambda=\frac{2h_m}{R(t)}(C_s-C_e),\qquad
\varepsilon_M=\overline C+\Lambda-2.55.
\end{equation}

问题一、三按1秒间隔、问题四按60秒间隔检查，均包含初始时刻和终点，并使用全部求解节点。最大平衡残差分别约为$4.44\times10^{-15}$、$4.88\times10^{-15}$、$7.11\times10^{-15}$。平均值和累计损失采用相同面通量构造，这些残差用于检查离散通量、积分权重及归一化水分收支的一致性。

\subsection{空间网格独立性}
采用160、320和640个区间逐级加密，在共同的物理位置与输出时刻比较。问题一覆盖前1800秒的全部逐秒点，问题二覆盖共同终端前的全过程逐秒点，问题三、四覆盖全部逐分钟点并加入共同终端。问题四比较实际域内位置及实际表面值，误差定义见式\eqref{eq:errors}。

\begin{equation}\label{eq:errors}
E_\theta=\max_{\mathcal S}|\theta_N-\theta_{2N}|,\quad
E_C=\max_{\mathcal S}|C_N-C_{2N}|,\quad
E_t=\frac{|t_{*,N}-t_{*,2N}|}{t_{*,2N}},
\end{equation}

其中$\mathcal S$为共同输出点集合。误差判据取$E_\theta<0.005$摄氏度、$E_C<5\times10^{-5}$；有终点事件的情景另要求$E_t<10^{-3}$。表\ref{tab:grid}列出两次加密的含水率差及最后一次加密的温度和事件差异，图\ref{fig:q3-convergence}展示场差随网格加密的变化。

\begin{table}[htbp]\centering\small
\caption{三网格比较；末三列为320与640区间之差}\label{tab:grid}
\begin{tabular}{crrrr}
\toprule
问题 & $E_C(160,320)$ & $E_C(320,640)$ & $E_\theta$/摄氏度 & $E_t$\\
\midrule
一 & $5.864\times10^{-5}$ & $1.468\times10^{-5}$ & $8.309\times10^{-6}$ & ---\\
二 & $5.130\times10^{-5}$ & $1.284\times10^{-5}$ & $1.444\times10^{-5}$ & ---\\
三 & $4.423\times10^{-5}$ & $1.106\times10^{-5}$ & $1.444\times10^{-5}$ & $2.434\times10^{-5}$\\
四 & $1.896\times10^{-4}$ & $4.742\times10^{-5}$ & $1.687\times10^{-5}$ & $1.511\times10^{-5}$\\
\bottomrule
\end{tabular}
\end{table}

\begin{figure}[htbp]\centering
\includegraphics[width=\linewidth]{figures/process_q3_convergence.pdf}
\caption{相邻网格的最大场差，$N$为细网格区间数；横线为含水率与温度误差判据，事件及时间精度结果见表\ref{tab:grid}、\ref{tab:time}。}\label{fig:q3-convergence}
\end{figure}

160至320区间的比较中，问题一、二、四的含水率差超过判据；320至640区间时，全部指标满足要求，故统一采用640区间。含水率差大体降为前一次的四分之一，呈现接近二阶的收敛趋势。

问题三、四在320与640区间下的事件时间分别相差约5.04秒和2.78秒，相对差均小于$10^{-3}$。整秒终点是当前离散模型按严格不等式给出的输出时刻，现有网格比较尚未提供亚秒级绝对误差保证。

\subsection{时间精度与终点稳定性}
在640区间网格上，将相对容差由$10^{-8}$收紧至$10^{-9}$，绝对容差相应收紧，结果见表\ref{tab:time}。四问的最大温度差不超过$1.843\times10^{-5}$摄氏度，最大含水率差不超过$1.589\times10^{-7}$；问题三、四的事件相对差分别为$1.749\times10^{-9}$和$3.510\times10^{-9}$。

\begin{table}[htbp]\centering\small
\caption{最终网格上两组时间容差的差异}\label{tab:time}
\begin{tabular}{crrr}
\toprule
问题 & 最大温度差/摄氏度 & 最大含水率差/(kg/kg) & 事件相对差\\
\midrule
一 & $8.635\times10^{-6}$ & $1.045\times10^{-7}$ & ---\\
二 & $1.476\times10^{-5}$ & $1.295\times10^{-7}$ & ---\\
三 & $1.476\times10^{-5}$ & $1.230\times10^{-7}$ & $1.749\times10^{-9}$\\
四 & $1.842\times10^{-5}$ & $1.589\times10^{-7}$ & $3.510\times10^{-9}$\\
\bottomrule
\end{tabular}
\end{table}

结合表\ref{tab:grid}和表\ref{tab:time}，含水率及终点时间对时间容差的敏感性小于对空间加密的敏感性。温度的两类差异均远低于0.005摄氏度的判据。上述比较支持当前数值设置下的场分布与终点稳定性，误差统计范围均为各检验定义的共同采样点集合。

\section{几何对照与工况灵敏度}
\subsection{分离物性改变与收缩影响}
为考察给定物性下的收缩影响，保持问题四的物性、初值、传递系数及环境延拓不变，仅将半径固定为2厘米。对照越阈时间为129.8449小时，长于收缩情景的51.0908小时。

记$\tau(\mathcal P,R)$为物性组$\mathcal P$和几何轨迹$R$下的越阈时间，$\mathcal P_3$、$\mathcal P_4$分别表示问题三、四的整组物性。现有计算构成$\tau(\mathcal P_3,R_0)\to\tau(\mathcal P_4,R_0)\to\tau(\mathcal P_4,R(t))$的比较路径。先改变物性使时间增加72.3721小时，再引入收缩使时间减少78.7541小时，净变化见式\eqref{eq:path-contrast}。

\begin{equation}\label{eq:path-contrast}
\begin{aligned}
\Delta\tau_{\rm total}
&=\underbrace{\tau(\mathcal P_4,R_0)-\tau(\mathcal P_3,R_0)}_{\text{固定几何下的整组物性差异}}
 +\underbrace{\tau(\mathcal P_4,R(t))-\tau(\mathcal P_4,R_0)}_{\text{问题四物性下的几何差异}}\\
&=72.3721-78.7541\approx-6.3819\ {\rm h}.
\end{aligned}
\end{equation}

差值由未舍入事件时间计算，展示时保留四位小数。沿这条路径，收缩引起的缩短超过整组物性改变引起的延长。该分解依赖所选比较路径，只描述整组物性与给定物性下几何变化的条件差异。

图\ref{fig:q4-counter}比较三种情景，式\eqref{eq:reductions}分别计算总体降幅和同物性下的收缩降幅。前者以问题三为基准，后者以问题四物性、固定半径情景为基准，两个百分比具有不同分母。

\begin{equation}\label{eq:reductions}
\begin{aligned}
\eta_{\rm total}
&=1-\frac{51.090831819}{57.472766796}=11.1043\%,\\
\eta_{\rm shrink}
&=1-\frac{51.090831819}{129.844890064}=60.6524\%.
\end{aligned}
\end{equation}

\begin{figure}[htbp]\centering
\includegraphics[width=\linewidth]{figures/result_q4_counterfactual.pdf}
\caption{三种情景的越阈时间：上、中两组几何相同，中、下两组物性相同。}\label{fig:q4-counter}
\end{figure}

问题四终点半径约为初始的0.6倍，半径平方约为初始的0.36倍，侧面积与体积之比约增至初始的1.667倍。这些变化分别影响内部扩散尺度和表面交换。实际半径与物性均随过程变化，完成时间由整个变化轨迹上的联立方程确定。

固定半径对照采用独立定义的恒定几何，即使计算超过72小时，也无需延拓收缩记录。该情景继续使用4小时后的恒定环境，用于比较给定物性下的几何影响。

\subsection{环境延拓和传质系数情景}
围绕两问的名义情景，分别采用末值延拓、末小时均值延拓、尾段温度加减0.5摄氏度、尾段有效边界含水率加减0.001以及传质系数加减10\%。环境延拓及温湿扰动从4小时后起作用，传质系数在全过程改变。每次只改变一项设置，越阈时间见表\ref{tab:sensitivity}。

\begin{table}[htbp]\centering\small
\caption{各单因素情景下的全域越阈时间，单位：小时}\label{tab:sensitivity}
\begin{tabular}{lrr}
\toprule
情景 & 问题三 & 问题四\\
\midrule
名义环境及传质系数 & 57.4728 & 51.0908\\
四小时后采用末值 & 57.1698 & 50.8245\\
四小时后采用末小时均值 & 57.4745 & 51.0923\\
尾段温度降低0.5摄氏度 & 58.3959 & 51.9004\\
尾段温度升高0.5摄氏度 & 56.5700 & 50.2989\\
尾段平衡含水率降低0.001 & 57.4532 & 51.0691\\
尾段平衡含水率升高0.001 & 57.4931 & 51.1137\\
全过程传质系数降低10\% & 58.2086 & 51.3820\\
全过程传质系数提高10\% & 56.8954 & 50.8700\\
\bottomrule
\end{tabular}
\end{table}

图\ref{fig:q4-sensitivity}给出相对各自名义时间的变化。尾段温度降低0.5摄氏度时，问题三、四的时间分别增加1.6062\%、1.5846\%；升高时分别减少1.5707\%、1.5500\%。末小时均值延拓与名义情景的差异约为0.003\%，末值延拓约缩短0.52\%。

\begin{figure}[htbp]\centering
\includegraphics[width=\linewidth]{figures/result_q4_sensitivity.pdf}
\caption{各单因素情景相对本问名义时间的变化，圆点与方块分别表示问题三和问题四。}\label{fig:q4-sensitivity}
\end{figure}

传质系数降低10\%时，问题三、四的时间分别增加1.2804\%、0.5699\%；提高10\%时，分别减少1.0046\%、0.4323\%。在所考察范围内，外部传质条件仍影响完成时间。

各扰动的单位、相对幅度和作用时间不同，结果用于描述所列情景的影响，不能据此给出统一的参数重要性排序或统计置信区间。

在全部已计算的成对情景中，问题四均早于问题三达到阈值，该时间排序在所考察的局部变化下保持一致。
